\chapter{Векторы состояния, унитарная эволюция и симметрия трансляций}
\providecommand{\ket}[1]{\left\lvert#1\right\rangle}
\providecommand{\bra}[1]{\left\langle#1\right\rvert}

После двухквартального перерыва полезно кратко повторить квантовую механику. Это не может занять весь квартал. Дальше речь пойдёт о применении квантовой механики к тем задачам, ради которых она первоначально и была разработана: к атомам, электронам и родственным частицам, а также к фотонам.

Сначала идут атомы. Как атомы могут идти раньше электронов, если электроны — части атомов? На первых порах нужна лишь такая механика электронов, которая позволяет записать основные уравнения атома водорода. Более подробное изучение электроноподобных частиц может последовать позже. Во всех этих приложениях симметрия будет постоянно возвращающимся инструментом: то, как симметрии представляются в квантовой механике и что они означают для системы.

\section{Состояния, наблюдаемые величины и ортогональность}

В классической механике состояние представляется точкой в фазовом пространстве. В квантовой механике оно представляется вектором в комплексном векторном пространстве. Это абстрактный вектор, а не стрелка в обычном трёхмерном пространстве. Векторы состояния можно складывать и умножать на комплексные числа, а пространство, которому они принадлежат, не обязано иметь какую-либо определённую конечную размерность.

В обозначениях Дирака вектор состояния записывают как кет,
\begin{equation}
\ket{\psi},
\end{equation}
а сопоставленный ему бра — как
\begin{equation}
\bra{\psi}.
\end{equation}
Бра и кеты связаны эрмитовым сопряжением. Если объединить бра и кет, получится скалярное произведение,
\begin{equation}
\langle\phi|\psi\rangle,
\end{equation}
которое, вообще говоря, является комплексным числом.

Наблюдаемая величина — это то, что можно измерить: координата, импульс, энергия или любая другая измеримая величина. В квантовой механике наблюдаемая величина представляется линейным эрмитовым оператором \(A\),
\begin{equation}
A^\dagger=A.
\end{equation}
Эрмитовость — операторный аналог вещественности. В частности, собственные значения эрмитова оператора вещественны.

Собственный вектор \(A\) удовлетворяет
\begin{equation}
A\ket{\alpha}=\alpha\ket{\alpha}.
\end{equation}
Возможные собственные значения — это возможные результаты измерения \(A\). Если, например, наблюдаемая величина может принимать только \(1\), \(3\) или \(7\), то эти числа и являются её собственными значениями. Её собственные векторы — это состояния, в которых результат определён, а не статистичен. В состоянии \(\ket{\alpha}\) измерение \(A\) с полной определённостью даёт \(\alpha\).

Скалярное произведение выражает важное логическое отношение между состояниями. Предположим, существует измерение, которое с полной определённостью различает \(\ket{\phi}\) и \(\ket{\psi}\). Тогда состояния ортогональны:
\begin{equation}
\langle\phi|\psi\rangle=0.
\end{equation}
Ортогональность означает, что состояния физически различимы без неоднозначности. Для нормированного состояния,
\begin{equation}
\langle\psi|\psi\rangle=1.
\end{equation}

\section{Частица в координатном представлении}

Сначала рассмотрим частицу, движущуюся вдоль прямой. Её координата \(x\) — наблюдаемая величина. Для каждого возможного положения \(x_0\) существует состояние \(\ket{x_0}\), в котором измерение координаты с полной определённостью даёт \(x_0\).

Состояния, соответствующие разным положениям, различимы и потому ортогональны:
\begin{equation}
x\neq x'
\quad\Longrightarrow\quad
\langle x'|x\rangle=0.
\end{equation}
Для произвольного состояния \(\ket{\psi}\) его перекрытие с координатным состоянием определяет волновую функцию:
\begin{equation}
\psi(x)=\langle x|\psi\rangle.
\end{equation}
Это координатное представление вектора состояния.

Волновая функция сама по себе не является вероятностью. Она может быть положительной, отрицательной, мнимой или вообще комплексной. Плотность вероятности для координаты равна
\begin{equation}
P(x)=\psi^*(x)\psi(x)=|\psi(x)|^2.
\end{equation}
Это произведение неотрицательно, как и должна быть плотность вероятности. Для нормированной волновой функции,
\begin{equation}
\int_{-\infty}^{\infty}dx\,|\psi(x)|^2=1.
\end{equation}

В координатном представлении оператор координаты действует умножением:
\begin{equation}
(\hat x\psi)(x)=x\psi(x).
\end{equation}
Его собственные функции — идеализированные функции, сосредоточенные в одной точке. Формально они описываются дельта-функциями Дирака: очень узкими пиками, центрированными в точке, соответствующей собственному значению координаты. Их точное определение как распределений для настоящего обзора не требуется.

В трёх измерениях координата имеет три компоненты, \(x\), \(y\) и \(z\), и координатное состояние представляет точку в трёхмерном пространстве. Два таких состояния ортогональны всякий раз, когда точки различаются хотя бы по одной координате. Обычно будет сохраняться одномерное обозначение, при подразумеваемом соответствующем трёхмерном обобщении.

\section{Импульс и состояния с определённым импульсом}

Каждой компоненте координаты соответствует компонента импульса. В координатном представлении оператор импульса вдоль оси \(x\) имеет вид
\begin{equation}
\hat p_x=-i\hbar\frac{d}{dx}.
\end{equation}
Таким образом, абстрактный эрмитов оператор можно также рассматривать конкретно как действие на волновые функции.

Состояние с определённым импульсом \(p_0\) удовлетворяет уравнению на собственные значения
\begin{equation}
\hat p_x\psi_{p_0}(x)=p_0\psi_{p_0}(x).
\end{equation}
Подставляя дифференциальную форму \(\hat p_x\), получаем
\begin{equation}
-i\hbar\frac{d\psi_{p_0}(x)}{dx}
=p_0\psi_{p_0}(x).
\end{equation}
Производная волновой функции пропорциональна самой волновой функции, поэтому решение имеет экспоненциальный вид:
\begin{equation}
\psi_{p_0}(x)\propto
\exp\left(\frac{ip_0x}{\hbar}\right).
\end{equation}

\begin{classroomqa}
\lecturetimestamp{00:24:26}
\audiencequestion{Выполняется ли это дифференциальное уравнение для любой волновой функции?}
\lecturerresponse{Нет. Это уравнение на собственные значения для состояния с определённым импульсом. Большинство волновых функций ему не удовлетворяют; оно выделяет собственные функции импульса.}
\end{classroomqa}

Различие между собственными состояниями координаты и импульса очень резкое. Собственная функция координаты сосредоточена в одной точке. Собственная функция импульса имеет
\begin{equation}
\left|\exp\left(\frac{ip_0x}{\hbar}\right)\right|^2=1.
\end{equation}
Поэтому её модуль постоянен на всей прямой. В этом идеализированном смысле определённый импульс означает полную неопределённость координаты; определённая координата точно так же означает полную неопределённость импульса. Это крайняя форма соотношения неопределённости координата--импульс.

\begin{classroomqa}
\lecturetimestamp{00:26:36}
\audiencequestion{Если известно, что импульс равен \(p_0\), почему всё равно остаётся распределение вероятности?}
\lecturerresponse{Показанная величина относится к измерению координаты, а не к ещё одному измерению импульса. Импульс и координата — разные наблюдаемые величины. Для идеализированного собственного состояния импульса координатная амплитуда имеет постоянный модуль на всей прямой. Амплитуды импульса относятся к импульсному представлению и связаны с координатной волновой функцией преобразованием Фурье.}
\end{classroomqa}

\begin{classroomqa}
\lecturetimestamp{00:28:02}
\audiencequestion{Как выглядит вероятность координаты, когда волновая функция является дельта-функцией Дирака?}
\lecturerresponse{Она сосредоточена в указанном положении. Идеализированный пик выражает полную определённость координаты.}
\end{classroomqa}

\editorialnote{Точные собственные состояния координаты и импульса являются обобщёнными состояниями. Дельта-функция Дирака и плоская волна не являются обычными квадратично интегрируемыми нормированными волновыми функциями на всей прямой, а буквальный квадрат дельта-функции не является обычным распределением. Сделанные выше утверждения касаются их идеализованной локализации и постоянного модуля.}

Волновую функцию можно формально рассматривать как вектор, имеющий по одной компоненте для каждого значения \(x\). Однако, в отличие от обычного конечного вектора-столбца, она имеет непрерывно бесконечное число компонент.

\begin{classroomqa}
\lecturetimestamp{00:28:53}
\audiencequestion{Представляются ли векторы состояния в виде векторов-столбцов?}
\lecturerresponse{Их можно представлять таким образом, но координатную волновую функцию неудобно записывать как обычный столбец, потому что \(x\) принимает непрерывно много значений. Её можно формально рассматривать как непрерывный вектор-столбец, а соответствующий бра — как непрерывный вектор-строку.}
\end{classroomqa}

\section{Эволюция во времени и унитарность}

Как состояние изменяется со временем? Эволюция во времени — это преобразование пространства состояний. Пусть \(U(t)\) обозначает эволюцию за прошедшее время \(t\). Тогда
\begin{equation}
U(t)\ket{\psi(t_1)}=\ket{\psi(t_1+t)}.
\end{equation}
В частности,
\begin{equation}
\ket{\psi(t)}=U(t)\ket{\psi(0)},
\qquad
U(0)=\mathbf 1.
\end{equation}

Условие на \(U(t)\) следует из сохранения информации. Если два начальных состояния различимы, эволюция во времени не должна стирать это различие. Ортогональные состояния должны оставаться ортогональными. Более общо, скалярное произведение любых двух состояний сохраняется:
\begin{equation}
\langle U(t)\phi|U(t)\psi\rangle
=\langle\phi|\psi\rangle.
\end{equation}
Поскольку
\begin{equation}
\langle U(t)\phi|U(t)\psi\rangle
=\langle\phi|U^\dagger(t)U(t)|\psi\rangle,
\end{equation}
сохранение всех скалярных произведений требует
\begin{equation}
U^\dagger(t)U(t)=\mathbf 1.
\end{equation}
Оператор, удовлетворяющий этому уравнению, является унитарным.

\begin{classroomqa}
\lecturetimestamp{00:36:36}
\audiencequestion{Является ли сохранение каждого скалярного произведения более сильным утверждением, чем простое сохранение ортогональности?}
\lecturerresponse{Для обсуждаемых линейных преобразований состояний сохранение ортогональности влечёт сохранение скалярных произведений. Обычные вращения дают привычную аналогию: они сохраняют ортогональность и скалярные произведения.}
\end{classroomqa}

\editorialnote{Строго говоря, одно лишь сохранение ортогональности допускает общее ненулевое масштабирование. Сохранение нормировки, уже требуемое для физической эволюции во времени, устраняет эту свободу и приводит к сохранению скалярных произведений.}

Унитарные операторы не следует путать с эрмитовыми наблюдаемыми величинами. Операторы координаты и импульса эрмитовы, но, вообще говоря, не унитарны. Собственные значения унитарного оператора — это фазы, то есть числа вида \(e^{i\theta}\), лежащие на единичной окружности.

\begin{classroomqa}
\lecturetimestamp{00:38:12}
\audiencequestion{Какой пример неунитарного преобразования можно привести?}
\lecturerresponse{Операторы координаты и импульса служат примерами: это эрмитовы наблюдаемые, но они не унитарны. Унитарный оператор имеет собственные значения единичного модуля, так что они являются фазовыми множителями.}
\end{classroomqa}

\section{Гамильтониан и Schr\"odinger уравнение}

Предполагается, что эволюция во времени непрерывна. Поэтому на очень малом интервале \(\epsilon\) оператор эволюции должен быть близок к единичному оператору:
\begin{equation}
U(\epsilon)=\mathbf 1+\epsilon G+\mathcal O(\epsilon^2).
\end{equation}
Используя унитарность,
\begin{align}
U^\dagger(\epsilon)U(\epsilon)
&=
\left(\mathbf 1+\epsilon G^\dagger\right)
\left(\mathbf 1+\epsilon G\right)
+\mathcal O(\epsilon^2)
\\
&=
\mathbf 1+\epsilon(G^\dagger+G)
+\mathcal O(\epsilon^2).
\end{align}
В первом порядке,
\begin{equation}
G^\dagger=-G.
\end{equation}
Следовательно, \(G\) антиэрмитов. Его можно записать в виде
\begin{equation}
G=-\frac{i}{\hbar}H,
\end{equation}
где \(H\) эрмитов:
\begin{equation}
H^\dagger=H.
\end{equation}
Поэтому бесконечно малый оператор эволюции во времени равен
\begin{equation}
U(\epsilon)
=
\mathbf 1-\frac{i\epsilon}{\hbar}H
+\mathcal O(\epsilon^2).
\end{equation}
Эрмитов оператор \(H\) — это гамильтониан.

\begin{classroomqa}
\lecturetimestamp{00:45:48}
\audiencequestion{Что позволяет пренебречь членом \(\epsilon^2\)?}
\lecturerresponse{Вычисление последовательно проводится только в первом порядке по \(\epsilon\). Вычисление второго порядка потребовало бы всех членов второго порядка, включая член, пропорциональный \(-\epsilon^2H^2/(2\hbar^2)\). Приближение заключено в усечённом выражении для \(U(\epsilon)\), а не в эрмитовости \(H\).}
\end{classroomqa}

Конечный промежуток времени можно составить из многих малых интервалов. Разобьём \(t\) на \(N\) шагов длины \(t/N\). Повторяя бесконечно малое преобразование, получаем
\begin{equation}
U(t)
=
\lim_{N\to\infty}
\left(
\mathbf 1-\frac{itH}{N\hbar}
\right)^N
=
\exp\left(-\frac{iHt}{\hbar}\right).
\end{equation}
Экспонента — это конечный аналог бесконечно малого преобразования.

Теперь проэволюционируем состояние от \(t\) до \(t+\epsilon\):
\begin{align}
\ket{\psi(t+\epsilon)}
&=
U(\epsilon)\ket{\psi(t)}
\\
&=
\left(
\mathbf 1-\frac{i\epsilon}{\hbar}H
\right)\ket{\psi(t)}
+\mathcal O(\epsilon^2).
\end{align}
Вычитая \(\ket{\psi(t)}\), деля на \(\epsilon\) и переходя к пределу, получаем
\begin{equation}
\frac{d}{dt}\ket{\psi(t)}
=
-\frac{i}{\hbar}H\ket{\psi(t)}.
\end{equation}
Эквивалентно,
\begin{equation}
i\hbar\frac{d}{dt}\ket{\psi(t)}
=
H\ket{\psi(t)}.
\end{equation}
Это зависящее от времени Schr\"odinger уравнение. В координатном представлении оно принимает вид
\begin{equation}
i\hbar\frac{\partial}{\partial t}\psi(x,t)
=
H\psi(x,t),
\end{equation}
где \(H\) действует как оператор на волновую функцию.

Не зависящее от времени Schr\"odinger уравнение есть уравнение на собственные значения гамильтониана:
\begin{equation}
H\ket{E}=E\ket{E}.
\end{equation}
Собственные значения \(E\) — это возможные измеряемые энергии, то есть энергетические уровни. Следовательно, один и тот же гамильтониан определяет и разрешённые энергии, и эволюцию состояния.

Если начальное состояние является собственным состоянием энергии, его эволюция имеет вид
\begin{equation}
\ket{E;t}
=
\exp\left(-\frac{iEt}{\hbar}\right)\ket{E}.
\end{equation}
Оно остаётся тем же собственным вектором энергии и приобретает только фазовый множитель.\footnote{Знак минус фиксируется предшествующим выводом \(U(t)=\exp(-iHt/\hbar)\). Более поздняя устная формулировка этой фазы имеет противоположный знак и несовместима с этим выводом.}

\begin{classroomqa}
\lecturetimestamp{01:01:05}
\audiencequestion{Как эволюционирует собственный вектор энергии согласно зависящему от времени Schr\"odinger уравнению?}
\lecturerresponse{Оставьте запись в виде кета. Поскольку \(H\ket{E}=E\ket{E}\), эволюция во времени умножает собственный вектор энергии на фазовый множитель, определяемый его энергией: \(\exp(-iEt/\hbar)\).}
\end{classroomqa}

\section{Что считается симметрией?}

Эволюция во времени — это одно унитарное преобразование в пространстве состояний. Преобразования симметрии образуют другой важный класс.

Вращательная симметрия играет центральную роль в атомной физике. В атоме водорода электрон движется в центральном кулоновском потенциале. Поворот системы координат не изменяет уравнения, и физический поворот атома по-прежнему оставляет атом водорода, подчиняющийся тем же уравнениям. Состояние при повороте может измениться, но динамический закон — нет.

Трансляция даёт ещё один пример. Изолированная система в пустом пространстве должна вести себя одинаково после того, как каждая её часть смещена на одну и ту же величину. Эта оговорка важна. Если присутствует печь и только система перемещается ближе к ней, система может вести себя иначе; печь нарушает эту трансляционную симметрию. Если же и система, и печь переносятся вместе, их исходное взаимное расположение восстанавливается.

Не всякая операция является симметрией. Сжатие кристалла каменной соли в 2 раза обычно изменяет его поведение. Сжатие — вполне определённое преобразование, но в общем случае не симметрия.

Следовательно, симметрия — это операция над состояниями, сохраняющая форму динамики. Обозначим её через \(V\). Она должна сохранять различимость: два ортогональных состояния после преобразования должны оставаться ортогональными. Соответственно,
\begin{equation}
V^\dagger V=\mathbf 1,
\end{equation}
так что \(V\) унитарен.

\section{Симметрии коммутируют с эволюцией во времени}

Предположим, что начальное состояние \(\ket{\psi_1}\) эволюционирует в \(\ket{\psi_2}\):
\begin{equation}
\ket{\psi_2}=U(t)\ket{\psi_1}.
\end{equation}
Применим предполагаемую симметрию \(V\) к обоим состояниям:
\begin{equation}
\ket{\psi_1'}=V\ket{\psi_1},
\qquad
\ket{\psi_2'}=V\ket{\psi_2}.
\end{equation}
Если \(V\) является симметрией динамики, то преобразованное начальное состояние должно эволюционировать в преобразованное конечное состояние:
\begin{equation}
\ket{\psi_2'}=U(t)\ket{\psi_1'}.
\end{equation}
Подставляя определения, получаем
\begin{equation}
VU(t)\ket{\psi_1}
=
U(t)V\ket{\psi_1}.
\end{equation}
Это должно выполняться для любого начального состояния. Следовательно,
\begin{equation}
VU(t)=U(t)V.
\end{equation}

Два эквивалентных пути можно изобразить коммутативной диаграммой.\footnote{Диаграмма восстанавливает аргумент с двумя путями, изложенный между 01:10:30 и 01:16:34.}
\begin{figure}
\centering
\begin{tikzpicture}[>=Stealth]
\node (a) at (0,0) {$\ket{\psi_1}$};
\node (b) at (4.2,0) {$\ket{\psi_2}=U(t)\ket{\psi_1}$};
\node (c) at (0,-2.0) {$V\ket{\psi_1}$};
\node (d) at (4.2,-2.0) {$V\ket{\psi_2}=U(t)V\ket{\psi_1}$};

\draw[->] (a) -- node[above] {$U(t)$} (b);
\draw[->] (c) -- node[below] {$U(t)$} (d);
\draw[->] (a) -- node[left] {$V$} (c);
\draw[->] (b) -- node[right] {$V$} (d);
\end{tikzpicture}
\caption{Операция симметрии и эволюция во времени дают одно и то же конечное состояние в любом порядке.}
\end{figure}

\begin{classroomqa}
\lecturetimestamp{01:16:38}
\audiencequestion{Должна ли симметрия коммутировать с эволюцией во времени при всех временах?}
\lecturerresponse{Да. Симметрия динамики должна удовлетворять условию коммутирования для полной эволюции во времени.}
\end{classroomqa}

\begin{classroomqa}
\lecturetimestamp{01:17:13}
\audiencequestion{Означает ли это, что можно повернуть, а потом подождать, или сначала подождать, а потом повернуть?}
\lecturerresponse{Именно. Если поворот является симметрией, оба порядка дают одно и то же преобразованное конечное состояние.}
\end{classroomqa}

Для малого интервала времени,
\begin{equation}
U(\epsilon)
=
\mathbf 1-\frac{i\epsilon}{\hbar}H
+\mathcal O(\epsilon^2).
\end{equation}
Условие коммутирования принимает вид
\begin{equation}
V\left(
\mathbf 1-\frac{i\epsilon}{\hbar}H
\right)
=
\left(
\mathbf 1-\frac{i\epsilon}{\hbar}H
\right)V
+\mathcal O(\epsilon^2).
\end{equation}
Члены с единичным оператором сокращаются, и остаётся
\begin{equation}
VH=HV,
\end{equation}
или
\begin{equation}
[V,H]=0.
\end{equation}
Таким образом, унитарная операция является симметрией динамики тогда и только тогда, когда она коммутирует с гамильтонианом.

\begin{classroomqa}
\lecturetimestamp{01:19:29}
\audiencequestion{Что требует, чтобы оператор симметрии \(V\) был унитарным?}
\lecturerresponse{Симметрия должна сохранять логические различия между состояниями. Она не должна переводить два взаимно исключающих, ортогональных состояния в состояния, которые уже нельзя различить. Сохранение этих отношений требует унитарного преобразования.}
\end{classroomqa}

Оператор без явной зависимости от времени, коммутирующий с \(H\), сохраняется. Это связывает симметрию с сохранением. В классической механике эта связь выражается теоремой Нётер. В данной квантово-механической форме ключевая проверка — коммутатор с гамильтонианом.

\section{Дискретные и непрерывные симметрии}

Симметрии могут быть дискретными или непрерывными. Зеркальное отражение дискретно: оно меняет местами левое и правое, причём между ними нет последовательности бесконечно малых отражений. Перестановка двух тождественных частиц — ещё одна дискретная операция. Переворачивание неразмеченного двустороннего предмета даёт пример того же типа.

Повороты непрерывны. Поворот на \(90^\circ\) можно составить из многих малых поворотов вокруг одной и той же оси. Пространственные трансляции непрерывны по той же причине: конечное смещение можно составить из многих очень малых смещений.

Эта делимость на малые преобразования придаёт непрерывным симметриям особую форму. Вблизи тождественного преобразования запишем
\begin{equation}
V(\epsilon)
=
\mathbf 1-i\epsilon G+\mathcal O(\epsilon^2),
\end{equation}
где \(\epsilon\) параметризует преобразование. Унитарность влечёт
\begin{equation}
G^\dagger=G.
\end{equation}
Эрмитов оператор \(G\) называется генератором симметрии.

Если преобразование является симметрией, то
\begin{equation}
[H,V(\epsilon)]=0.
\end{equation}
Разложение по первому порядку даёт
\begin{equation}
[H,G]=0.
\end{equation}
Следовательно, непрерывные симметрии порождаются эрмитовыми операторами, коммутирующими с гамильтонианом. Эти генераторы являются сохраняющимися величинами.

Слово ``порождаются'' здесь имеет конкретный смысл. Конечное преобразование строится многократным применением бесконечно малого, точно так же как конечная эволюция во времени строится из многих коротких временных шагов.

\section{Трансляционная симметрия и импульс}

Рассмотрим активную трансляцию волновой функции вправо на расстояние \(\epsilon\). Если исходная волновая функция сосредоточена около \(x_0\), то сдвинутая должна быть сосредоточена около \(x_0+\epsilon\). Поэтому её значение в координате \(x\) равно значению исходной волновой функции в \(x-\epsilon\):
\begin{equation}
(T(\epsilon)\psi)(x)=\psi(x-\epsilon).
\end{equation}
Для малого смещения,
\begin{equation}
\psi(x-\epsilon)
=
\psi(x)-\epsilon\frac{d\psi(x)}{dx}
+\mathcal O(\epsilon^2).
\end{equation}
Следовательно,
\begin{equation}
T(\epsilon)
=
\mathbf 1-\epsilon\frac{d}{dx}
+\mathcal O(\epsilon^2).
\end{equation}
Используя
\begin{equation}
\hat p_x=-i\hbar\frac{d}{dx},
\qquad
\frac{d}{dx}=\frac{i}{\hbar}\hat p_x,
\end{equation}
бесконечно малая трансляция принимает вид
\begin{equation}
T(\epsilon)
=
\mathbf 1-\frac{i\epsilon}{\hbar}\hat p_x
+\mathcal O(\epsilon^2).
\end{equation}
Соответствующая конечная трансляция равна
\begin{equation}
T(\epsilon)
=
\exp\left(-\frac{i\epsilon}{\hbar}\hat p_x\right).
\end{equation}
При соглашении \(V(\epsilon)=\mathbf 1-i\epsilon G+\cdots\) для параметра трансляции с размерностью получаем
\begin{equation}
G_{\mathrm{trans}}=\frac{\hat p_x}{\hbar}.
\end{equation}
Эквивалентно, импульс называют генератором пространственных трансляций, когда множитель \(\hbar\) явно выписан в экспоненте.

\begin{classroomqa}
\lecturetimestamp{01:34:08}
\audiencequestion{Не потерялся ли знак минус в выводе для трансляции?}
\lecturerresponse{Здесь вопрос в направлении, приписываемом сдвинутой волновой функции. Активный сдвиг вправо представляется как \(\psi(x-\epsilon)\), а не как \(\psi(x+\epsilon)\). С этой поправкой получается \(T(\epsilon)=\mathbf 1-i\epsilon\hat p_x/\hbar+\cdots\).}
\end{classroomqa}

Сохраняется ли импульс, зависит от гамильтониана. Для свободной частицы, движущейся в одном измерении,
\begin{equation}
H=\frac{\hat p_x^2}{2m}.
\end{equation}
Поскольку любой оператор коммутирует с любой функцией от самого себя,
\begin{equation}
[H,\hat p_x]=0.
\end{equation}
Следовательно, гамильтониан свободной частицы инвариантен относительно трансляций, и \(\hat p_x\) сохраняется. Физически частица ведёт себя одинаково, где бы ни было выбрано её начальное положение. Сдвиг начального условия просто сдвигает всё движение.

В трёх измерениях,
\begin{equation}
H
=
\frac{\hat p_x^2+\hat p_y^2+\hat p_z^2}{2m}.
\end{equation}
Каждая компонента импульса коммутирует с \(H\):
\begin{equation}
[H,\hat p_x]=[H,\hat p_y]=[H,\hat p_z]=0.
\end{equation}
В каждом пространственном направлении имеется отдельная симметрия относительно трансляций. Импульс одновременно является сохраняющейся величиной и генератором, связанным с этими трансляциями.

Следующий случай — вращательная симметрия. Повороты устроены сложнее, поскольку повороты вокруг разных осей, вообще говоря, не коммутируют: поворот вокруг \(x\), за которым следует поворот вокруг \(y\), не совпадает с обратным порядком. Это приводит к некоммутирующим генераторам симметрии и к группо-теоретической структуре поворотов.
